\section{From the action inequality to spectral moments}
\label{sec:spectral}

We apply Theorem~\ref{thm:action} first to ordered Dirichlet
eigenfunctions.  We then pass to the line using finite-dimensional trial
spaces.  This last step treats all negative eigenvalues, without assuming
in advance that their moment is finite.

Throughout this section, $0<\sigma<1$, $p=1+\sigma$,
$q=1+1/\sigma$, and $\gamma=\sigma+1/2$.  Write
\begin{equation}\label{spectral:constant}
 A_\sigma=\frac{\sigma^\sigma}{(1+\sigma)^{1+\sigma}},
 \qquad
 \ell_\sigma=\frac{A_\sigma}{\mathrm B(1/2,1+\sigma)},
\end{equation}
where $\mathrm B$ denotes the beta function.

\subsection{Quadratic forms and negative spectrum}

For $0\leq W\in L^p(\R)$, define
\[
 h_{-W}[v]=\int_\R |v'|^2-\int_\R W|v|^2,
 \qquad v\in H^1(\R).
\]
We record the facts needed to define its operator and to exhaust its
negative eigenvalues. We use the association of a densely defined closed
semibounded form with its self-adjoint operator, and the spectral
description of the operator and its form domain, in
\cite[Theorem~2.13, Theorems~3.7--3.8 and (3.50)--(3.51)]{Teschl2009}.
All theorem numbers here refer to the 2009 first edition.

\newpage
\begin{lemma}\label{spectral:forms}
Let $p>1$ and $0\leq W\in L^p(\R)$.  The form $h_{-W}$ is closed
and bounded below on $H^1(\R)$.  For its associated self-adjoint operator
$H_{-W}$, every spectral projection onto $(-\infty,-\alpha]$, with
$\alpha>0$, has finite rank.  Thus its negative spectrum consists of
eigenvalues of finite multiplicity, with possible accumulation only at
zero.  On a bounded interval $I$, the corresponding Dirichlet form on
$H^1_0(I)$ is closed and bounded below and has compact resolvent.
\end{lemma}

\begin{proof}
Put $p'=p/(p-1)$.  The one-dimensional Sobolev inequality and
interpolation give
\[
 \|v\|_{2p'}^2
 \leq \|v\|_\infty^{2/p}\|v\|_2^{2-2/p}
 \leq 2^{1/p}\|v'\|_2^{1/p}\|v\|_2^{2-1/p}.
\]
Consequently H\"older's and Young's inequalities imply that, for every
$\eta>0$, there is $b_{\eta,W}<\infty$ such that
\begin{equation}\label{spectral:form-bound}
 \int_\R W|v|^2
 \leq \eta\|v'\|_2^2+b_{\eta,W}\|v\|_2^2.
\end{equation}
The KLMN Theorem \cite[Theorem~6.24]{Teschl2009} now gives closure and semiboundedness, with
form domain exactly $H^1(\R)$.

Multiplication by $\sqrt W$ is compact from $H^1(\R)$ to $L^2(\R)$.
Indeed, choose nonnegative bounded, compactly supported $W_n$ converging
to $W$ in $L^p$.  Multiplication by $\sqrt{W_n}$ is compact by local
Rellich compactness, and
\[
 \| (\sqrt W-\sqrt{W_n})v\|_2^2
 \leq \|W-W_n\|_p\|v\|_{2p'}^2
 \leq C_p\|W-W_n\|_p\|v\|_{H^1}^2.
\]
Thus these multiplication operators converge in operator norm.

If the spectral subspace of $H_{-W}$ for $(-\infty,-\alpha]$ were
infinite-dimensional, it would contain an orthonormal sequence $v_n$.
Semiboundedness places this entire spectral subspace in the form domain.
The lower bound obtained from \eqref{spectral:form-bound} with
$\eta=1/2$, together with $h_{-W}[v_n]\leq-\alpha$, bounds this
sequence in $H^1$.  After passing to a subsequence, it converges weakly
to zero in $H^1$, because it is orthonormal in $L^2$.  Compactness then
gives $\sqrt Wv_n\to0$ in $L^2$, contradicting
\[
 -\alpha\geq h_{-W}[v_n]\geq-\|\sqrt Wv_n\|_2^2.
\]
This proves the negative-spectrum assertion.  On a bounded interval,
the same form estimate follows by zero extension.  The form norm is
equivalent to the $H^1_0$ norm, whose embedding into $L^2(I)$ is compact;
hence the Dirichlet operator has compact resolvent.
\end{proof}

\subsection{The Dirichlet estimate}

The order of the eigenfunctions is what makes the triangular
coefficient restriction in Theorem~\ref{thm:action} useful.

\begin{proposition}\label{spectral:dirichlet}
Let $I$ be a bounded interval and $0\leq W\in L^p(I)$.  If
$-k_1^2,\ldots,-k_N^2$ are the first $N$ negative Dirichlet eigenvalues
of $-d^2/dx^2-W$, ordered so that $k_1\geq\cdots\geq k_N>0$, then
\begin{equation}\label{spectral:dirichlet-bound}
 \sum_{j=1}^N k_j^{2\gamma}
 \leq \ell_\sigma\int_I W^p.
\end{equation}
\end{proposition}

\begin{proof}
Choose real orthonormal eigenfunctions $u_1,\ldots,u_N$; these exist
because the Dirichlet form is real.  Put $u=(u_1,\ldots,u_N)^{\mathsf T}$
and $K=\operatorname{diag}(k_1,\ldots,k_N)$.  For every real lower
triangular matrix $C$, write $a=Cu$.  Its $i$th component lies in the
span of $u_1,\ldots,u_i$, and the eigenfunction equations in form sense
give
\[
 \int_I \bigl(|a_i'|^2-W|a_i|^2+k_i^2|a_i|^2\bigr)
 =\sum_{j\leq i} C_{ij}^2(k_i^2-k_j^2)\leq0.
\]
Summing over $i$ and putting $m=|a|^2$, we obtain
\[
 \mathcal E_K(a)\leq\int_I(Wm-m^q)
 \leq A_\sigma\int_I W^{1+\sigma}.
\]
The last inequality follows by maximizing $Wm-m^q$ over $m\geq0$;
the maximizer is $(W/q)^\sigma$.

Theorem~\ref{thm:action} therefore yields
\[
 \sum_{j=1}^N\int_{-k_j}^{k_j}(k_j^2-s^2)^\sigma\,ds
 \leq A_\sigma\int_I W^{1+\sigma}.
\]
Substituting $s=k_jt$ in each integral gives
$\mathrm B(1/2,1+\sigma)k_j^{2\sigma+1}$, proving
\eqref{spectral:dirichlet-bound}.
\end{proof}

The beta--gamma identity identifies the constant without any limiting
argument in $\sigma$:
\begin{align}\label{spectral:constant-identification}
 \ell_\sigma
 &=\left(\frac{\sigma}{1+\sigma}\right)^\sigma
   \frac{\Gamma(\sigma+3/2)}{\sqrt\pi\,\Gamma(\sigma+2)}
 \notag\\
 &=2\left(\frac{\sigma}{1+\sigma}\right)^\sigma
   L_{\gamma,1}^{\mathrm{cl}}.
\end{align}

\subsection{Exhaustion of the line}

Let $0\leq W\in L^p(\R)$.  If there are no negative eigenvalues, the
trace inequality is immediate.  Otherwise fix any finite number $N$
of them, counted with multiplicity and ordered increasingly:
\[
 \lambda_1\leq\cdots\leq\lambda_N<0.
\]
Let $\phi_1,\ldots,\phi_N\in H^1(\R)$ be corresponding orthonormal
eigenfunctions.  Choose smooth cutoffs $\chi_R$ equal to one on
$[-R,R]$, supported in $I_R=(-2R,2R)$, and satisfying
$0\leq\chi_R\leq1$ and $\|\chi_R'\|_\infty\leq c/R$.  Then
\[
 v_{j,R}=\chi_R\phi_j\longrightarrow\phi_j
 \quad\hbox{in }H^1(\R).
\]
H\"older's inequality and the embedding $H^1\subset L^{2p'}$ show
that the potential form is continuous under this convergence.

For clarity, let $G_R$ and $Q_R$ be the Gram and form matrices of the
$N$ functions $v_{j,R}$.  Then
\[
 G_R\longrightarrow I_N,
 \qquad
 Q_R\longrightarrow\Lambda
 :=\operatorname{diag}(\lambda_1,\ldots,\lambda_N)
\]
in matrix norm.  In particular the functions are linearly independent
for large $R$.  If $g_R=\|G_R-I_N\|_{\mathrm{op}}<1$ and
$e_R=\|Q_R-\Lambda\|_{\mathrm{op}}$, their Rayleigh quotients differ
uniformly from the limiting coefficient-space quotients by at most
\[
 \eta_R=\frac{e_R+\|\Lambda\|_{\mathrm{op}}g_R}{1-g_R}
 \longrightarrow0.
\]
Applying the form-domain min--max principle
\cite[Theorem~4.10 and the following form-domain extension]{Teschl2009}
to the span of
$v_{1,R},\ldots,v_{j,R}$ gives, for every $j\leq N$,
\[
 \lambda_j(I_R)\leq\lambda_j+\eta_R,
\]
where $\lambda_j(I_R)$ is the $j$th Dirichlet eigenvalue on $I_R$.
All these levels are negative for sufficiently large $R$.
Proposition~\ref{spectral:dirichlet} now implies
\[
 \sum_{j=1}^N(-\lambda_j-\eta_R)^\gamma
 \leq \sum_{j=1}^N|\lambda_j(I_R)|^\gamma
 \leq \ell_\sigma\int_{I_R}W^p
 \leq \ell_\sigma\int_\R W^p.
\]
Letting $R\to\infty$ proves the estimate for every finite partial
sum.  Taking their supremum gives
\begin{equation}\label{spectral:line-bound}
 \Tr(H_{-W})_-^\gamma
 \leq \ell_\sigma\int_\R W^p.
\end{equation}
Thus the full moment is finite even when there are infinitely many
negative eigenvalues.  No approximation of $W$ by smoother potentials
is needed.

\subsection{Signed potentials}

To include unbounded positive parts, we specify the form domain in the
signed-potential statement.

\begin{corollary}\label{spectral:signed}
Let $V$ be real measurable with $V_-\in L^p(\R)$, and suppose
\[
 \mathcal D_V
 =\left\{v\in H^1(\R):\int_\R V_+|v|^2<\infty\right\}
\]
is dense in $L^2(\R)$.  This holds in particular if
$V_+\in L^1_{\mathrm{loc}}(\R)$.  The form
\[
 h_V[v]=\int_\R|v'|^2+\int_\R V_+|v|^2-\int_\R V_-|v|^2,
 \qquad v\in\mathcal D_V,
\]
is closed and bounded below.  Its associated operator $H_V$ has
discrete negative spectrum and satisfies
\[
 \Tr(H_V)_-^\gamma\leq\ell_\sigma\int_\R V_-^p.
\]
\end{corollary}

\begin{proof}
The positive form $\int|v'|^2+\int V_+|v|^2$ is closed on
$\mathcal D_V$, by completeness of $H^1$ and of the weighted $L^2$
space.  Estimate~\eqref{spectral:form-bound}, with $W=V_-$, proves
closure and semiboundedness after subtraction of the negative part.
If $V_+\in L^1_{\mathrm{loc}}$, then
$C_c^\infty(\R)\subset\mathcal D_V$, proving the asserted density.
The compactness argument in Lemma~\ref{spectral:forms} still applies.
On a negative spectral subspace bounded away from zero, the same relative
form bound controls the $H^1$ norm of unit vectors, since the positive
potential term can be discarded. Compactness of multiplication by
$\sqrt{V_-}$ then excludes an infinite orthonormal sequence in that
subspace: the integral against $V_-$ would tend to zero along a subsequence,
whereas its negative form energy stays bounded away from zero. Thus every
such spectral projection has finite rank.

For $v\in\mathcal D_V\subset H^1(\R)$, one has
$h_V[v]\geq h_{-V_-}[v]$.  The min--max principle, with this inclusion
of form domains, bounds each negative eigenvalue of $H_V$ below by
the corresponding negative eigenvalue of $H_{-V_-}$; in particular,
if $H_V$ has $j$ negative eigenvalues, so does $H_{-V_-}$. Summation and
\eqref{spectral:line-bound} prove the result.
\end{proof}

\subsection{Sharpness}

Put $r=1/\sigma>1$ and consider
\[
 W(x)=(r+1)\sech^2(rx),
 \qquad
 \psi(x)=\sech^{1/r}(rx).
\]
Direct differentiation gives
\[
 \frac{\psi'}{\psi}=-\tanh(rx),
 \qquad
 \frac{\psi''}{\psi}=1-(r+1)\sech^2(rx).
\]
The function $\psi$ and its first two derivatives decay exponentially.
Thus $\psi\in H^2(\R)$, the operator domain for this bounded $W$,
and $H_{-W}\psi=-\psi$.
Substitution $t=rx$, followed by $z=\tanh t$, gives
\[
 \int_\R W^p
 =\frac{(r+1)^p}{r}\,\mathrm B(1/2,p),
 \qquad
 \left(\int_\R W^p\right)^{-1}=\ell_\sigma.
\]
The eigenvalue $-1$ gives $\Tr(H_{-W})_-^\gamma\geq1$; the upper
bound \eqref{spectral:line-bound} gives the reverse inequality.
This completes the proof of Theorem~\ref{thm:main}.
Hence this potential attains equality, and the constant in
\eqref{spectral:constant-identification} is optimal.  This establishes
$L_{\gamma,1}=L_{\gamma,1}^{(1)}$ throughout $1/2<\gamma<3/2$.
